% ============================================================================
%  TRACE-ESUS / Causal-ESUS -- formal statement of the latent-endotyping
%  data-generating process.
%
%  Source of truth: r21_latent_endotyping_experiment.py
%      SimulationConfig            lines  78--101
%      simulate_two_mechanism_cohort lines 232--274
%  Numerical check: outputs_latent_endotyping/recovery_summary.csv (oracle row)
%
%  Requires: amsmath, amssymb, bm  (bm optional; swap \bm for \mathbf if absent)
% ============================================================================

\subsection{Data-generating process}
\label{sec:generator}

Let $i = 1,\dots,N$ index patients. For each patient we define a binary renal
status $R_i \in \{0,1\}$ (observed), a categorical stroke mechanism
$Z_i \in \{A, C\}$ denoting the atrial-thromboembolic ($A$) and competing ($C$)
mechanisms (\emph{latent, never passed to any fitting routine}), and a vector of
three continuous biomarkers $\bm{B}_i \in \mathbb{R}^{3}$ (observed), indexed
$j \in \{0,1,2\}$ corresponding to an NT-proBNP-like marker, atrial electrical
evidence, and competing-mechanism evidence respectively.

\paragraph{Structural equations.}
The generative model is the structural causal model

\begin{align}
  R_i        &\sim \mathrm{Bernoulli}(\pi_R), \label{eq:gen-renal}\\[2pt]
  Z_i \mid R_i &\sim \mathrm{Categorical}\!\big(p(R_i),\, 1 - p(R_i)\big),
                 \qquad p(0) = p(1) = \tfrac{1}{2}, \label{eq:gen-mech}\\[2pt]
  \bm{B}_i   &= \bm{\mu}_{Z_i} \;+\; R_i\,\bm{\gamma} \;+\; \bm{\varepsilon}_i,
                 \qquad
                 \bm{\varepsilon}_i \sim \mathcal{N}(\bm{0}, \bm{\Sigma}),
                 \label{eq:gen-biomarkers}
\end{align}

with $\bm{\Sigma} = \operatorname{diag}(\sigma_0^2, \sigma_1^2, \sigma_2^2)$.
Componentwise, Eq.~\eqref{eq:gen-biomarkers} reads

\begin{equation}
  B_{ij} \;=\; \mu_{Z_i,\,j} \;+\; R_i\,\gamma_j \;+\; \varepsilon_{ij},
  \qquad
  \varepsilon_{ij} \sim \mathcal{N}(0, \sigma_j^2)
  \;\;\text{independently over } j .
  \label{eq:gen-component}
\end{equation}

\paragraph{Parameter values.}
All effects are expressed in units of within-biomarker residual standard
deviation; setting $\sigma_j \equiv 1$ makes this normalisation exact.

\begin{equation}
  \bm{\mu}_A =
  \begin{pmatrix} 1.25 \\ 1.00 \\ 0.00 \end{pmatrix}\!,
  \qquad
  \bm{\mu}_C =
  \begin{pmatrix} 0.00 \\ 0.00 \\ 1.00 \end{pmatrix}\!,
  \qquad
  \bm{\gamma} =
  \begin{pmatrix} \gamma_0 \\ 0 \\ 0 \end{pmatrix}\!,
  \qquad
  \bm{\Sigma} = \bm{I}_3,
  \label{eq:gen-params}
\end{equation}

with $\pi_R = 0.30$ and the renal distortion magnitude swept over
$\gamma_0 \in \{0.00,\, 0.50,\, 1.00,\, 1.50\}$, labelled
\textsc{None}, \textsc{Weak}, \textsc{Moderate}, \textsc{Strong}.

Two features of Eq.~\eqref{eq:gen-params} are load-bearing. First, the zero
entries of $\bm{\mu}_A$ and $\bm{\mu}_C$ encode \emph{prohibited} mechanism
$\to$ biomarker edges and are what render the markers diagnostic. Second,
$\bm{\gamma}$ is supported on a single coordinate: renal dysfunction distorts
the NT-proBNP-like marker and is structurally forbidden from touching the
remaining two. The Boolean support of $\bm{\gamma}$,
$\operatorname{supp}(\bm{\gamma}) = (\text{true}, \text{false}, \text{false})$,
is exactly the \texttt{renal\_path\_mask} supplied to the biologically
constrained model.

\paragraph{Graph and factorisation.}
The DAG is $Z \rightarrow \bm{B} \leftarrow R$, and the joint factorises as

\begin{equation}
  p(R, Z, \bm{B}) \;=\; p(R)\; p(Z \mid R)\; p(\bm{B} \mid Z, R)
                 \;=\; p(R)\; p(Z)\; p(\bm{B} \mid Z, R),
  \label{eq:gen-factorisation}
\end{equation}

the second equality following from $p(0) = p(1) = \tfrac12$ in
Eq.~\eqref{eq:gen-mech}, i.e.\ from the \emph{absence} of an $R \rightarrow Z$
edge. Consequently $Z \perp\!\!\!\perp R$ marginally. Renal status is therefore
\textbf{not} a confounder of the mechanism--biomarker relation; it is a
non-causal distortion of the measurement channel, and we refer to it throughout
as \emph{renal biomarker distortion}. Crucially, however,
$Z \not\perp\!\!\!\perp R \mid \bm{B}$: conditioning on the collider $\bm{B}$
renders renal status informative about mechanism, which is the formal basis for
the gain of the renal-adjusted and biologically constrained models over the
pooled associative baseline.

\paragraph{Oracle posterior.}
The data-generating posterior admits an exact reduction. Writing
$\tilde{\bm{B}}_i = \bm{B}_i - R_i \bm{\gamma}$,

\begin{equation}
  \Pr(Z_i = A \mid \bm{B}_i, R_i)
  = \frac{\pi_A\, \mathcal{N}(\bm{B}_i;\, \bm{\mu}_A + R_i\bm{\gamma},\, \bm{\Sigma})}
         {\displaystyle\sum_{z \in \{A,C\}}
           \pi_z\, \mathcal{N}(\bm{B}_i;\, \bm{\mu}_z + R_i\bm{\gamma},\, \bm{\Sigma})}
  = \Pr(Z_i = A \mid \tilde{\bm{B}}_i),
  \label{eq:oracle}
\end{equation}

since every density depends on $\bm{B}_i$ only through
$\bm{B}_i - \bm{\mu}_z - R_i\bm{\gamma} = \tilde{\bm{B}}_i - \bm{\mu}_z$. The
renal contribution is a common additive offset to both hypotheses and is thus
\emph{exactly removable by subtraction} rather than requiring adjustment. The
log-odds are linear in the corrected biomarkers,

\begin{equation}
  \log \frac{\Pr(Z_i = A \mid \bm{B}_i, R_i)}{\Pr(Z_i = C \mid \bm{B}_i, R_i)}
  = \log\frac{\pi_A}{\pi_C}
    + (\bm{\mu}_A - \bm{\mu}_C)^{\!\top} \bm{\Sigma}^{-1}
      \Big( \tilde{\bm{B}}_i - \tfrac{1}{2}(\bm{\mu}_A + \bm{\mu}_C) \Big),
  \label{eq:logodds}
\end{equation}

which for the parameter values of Eq.~\eqref{eq:gen-params} and
$\pi_A = \pi_C = \tfrac12$ evaluates to the discriminant

\begin{equation}
  s(\bm{B}_i, R_i)
  = 1.25\,\tilde{B}_{i0} + 1.00\,\tilde{B}_{i1} - 1.00\,\tilde{B}_{i2} - 0.78125,
  \qquad \tilde{B}_{i0} = B_{i0} - R_i \gamma_0 ,
  \label{eq:discriminant}
\end{equation}

with the atrial mechanism assigned when $s > 0$.

\paragraph{Irreducible error.}
Because the two class-conditional distributions are Gaussian with common
covariance and equal priors, the Bayes error is available in closed form as
$\Phi(-d/2)$, where $d$ is the Mahalanobis separation

\begin{equation}
  d = \big\lVert \bm{\mu}_A - \bm{\mu}_C \big\rVert_{\bm{\Sigma}}
    = \sqrt{1.25^2 + 1.00^2 + (-1.00)^2}
    = \sqrt{3.5625}
    = 1.8875 .
  \label{eq:separation}
\end{equation}

Hence the ceiling on top-1 mechanism accuracy is

\begin{equation}
  1 - \Phi(-d/2) = 1 - \Phi(-0.94375) = 0.82735 .
  \label{eq:bayes}
\end{equation}

Note that Eqs.~\eqref{eq:oracle} and~\eqref{eq:bayes} are independent of
$\gamma_0$: the oracle is by construction immune to renal distortion. The
simulated oracle reproduces this constant to Monte Carlo precision across all
four distortion levels ($0.8275$, $0.8265$, $0.8285$, $0.8270$ at
$\gamma_0 = 0.00, 0.50, 1.00, 1.50$; $500$ repeats per level). Any accuracy lost
by a fitted method as $\gamma_0$ increases is therefore attributable to
modelling failure rather than to destruction of information in the data.
